    \section{Basis}
        \subsection{Definition of basis}
            \hskip \parindent
            \par Let $V$ be a vector space over a field $\mathbb{F}$. A basis for $V$ is a l.i. set $B$ s.t.
            \begin{equation}
                \operatorname{span}\{B\} = V.
            \end{equation}
        \subsection{Example of basis (V=$\mathbb{R}^2$)}
            \subsubsection{$\{(1,0),(0,1),(1,1)\}$}
                \hskip \parindent
                \par $B=\{(1,0),(0,1),(1,1)\}$ is not a basis for $\mathbb{R}^2$. Because $B$ is not l.i. .
            \subsubsection{$\{(1,1),(1,-1)\}$}
                \hskip \parindent
                \par $S=\{(1,1),(1,-1)\}$ is a l.i. set and $\operatorname{span}S=\mathbb{R}^2$, so $S$ is a basis for $\mathbb{R}^2$.
                \par (x,y) = $\frac{x+y}{2}$(1,1) + $\frac{x-y}{2}$(1,-1)
        \subsection{Extension: unique representation of vectors in terms of basis}
            \hskip \parindent
            \par A set of vectors $v_1,v_2,\cdots,v_n$ is a basis for the vector space $V$ iff any $v \in V$ can be written uniquely as a l.c. $v=\alpha_1 v_1 + \alpha_2 v_2 + \cdots + \alpha_n v_n$.
            \par That is for each $v \in V$,
            \begin{equation}
                \exists! \alpha_1,\alpha_2,\cdots,\alpha_n \in \mathbb{F} \text{ s.t. } v=\alpha_1 v_1 + \alpha_2 v_2 + \cdots + \alpha_n v_n
            \end{equation}
        \subsection{Proposition: characterization of basis for $\mathbb{F}^m$\label{sec:characterizationofbasis}}
            \hskip \parindent
            \par A set of $n$ vectors $v_1,v_2,\cdots,v_n \in \mathbb{F}^m$ is a basis for $\mathbb{F}^m$ iff $n=m$ and the rank of $A = [v_1|v_2|\cdots|v_n]$ is $n$.
            \par Proof see section "Proofs".
        \subsection{Theorem: characterization of basis in general\label{sec:theorem_characterization_of_basis}}
            \hskip \parindent
            \par Let V be a f.d.v.s. over a field $\mathbb{F}$ with $dim_{\mathbb{F}}V=n$. Let $\{v_1,v_2,\cdots,v_n\}$ be a set of vectors in V. Then the following conditions are equivalent:
            \begin{itemize}
                \item [(1)]: $\{v_1,v_2,\cdots,v_n\}$ is a basis for V.
                \item [(2)]: V=$\operatorname{span}\{v_1,v_2,\cdots,v_n\}$
                \item [(3)]: $\{v_1,v_2,\cdots,v_n\}$ is a l.i. set.
            \end{itemize}
            \par The proof see section "Proofs".
        \subsection{Example of the characterization of basis}
            \hskip \parindent
            \par The vectors $(1,-1)$ and $(1,1)$ form a basis for $\mathbb{R}^2$. $A = \begin{bmatrix}
                1 & 1 \\
                -1 & 1
            \end{bmatrix}$, $\text{rank}(A)=2$.
            \par Therefore, $\{(1,-1),(1,1)\}$ is a basis for $\mathbb{R}^2$.
        \subsection{Extension: span of subsets of a spanning set}
            \hskip \parindent
            \par If $w_1,w_2,\cdots,w_k$ are vectors in span$\{v_1,v_2,\cdots,v_n\}$, then span$\{w_1,w_2,\cdots,w_k\}$ $\subseteq$ span$\{v_1,v_2,\cdots,v_n\}$.
        \subsection{Extension: add a vector already in the span}
            \hskip \parindent
            \par Let $\{v_1,v_2,\cdots,v_r\}$ be a set of vectors such that $v_r \in \operatorname{span}\{v_1,v_2,\cdots,v_{r-1}\}$.
            \par Then
            \[
                \operatorname{span}\{v_1,\ldots,v_{r-1},v_r\} = \operatorname{span}\{v_1,\ldots,v_{r-1}\},
            \]
            \par i.e. adding $v_r$ does not enlarge the span.
        \subsection{Theorem: basis subset of a finite generating set\label{sec:theorem_basis_subset}}
            \hskip \parindent
            \par Let $\{v_1,v_2,\cdots,v_n\}$ be a finite set of non-zero vectors in a vector space $V$. Then $\{v_1,v_2,\cdots,v_n\}$ contains a basis for the subspace $\operatorname{span}\{v_1,v_2,\cdots,v_n\}$.
            \par Proof see section "Proofs".
        \subsection{Proposition: existence of basis\label{sec:existenceofbasis}}
            \hskip \parindent
            \par Let $v_1,v_2,\cdots,v_n$ be vectors in $\mathbb{F}^m$ and $A=[v_1|v_2|\cdots|v_n]\in M_{m\times n}(\mathbb{F})$.
            \par Let $R$ be a row-reduced echelon form that is row-equivalent to $A$.
            \par Let $i_1<i_2<\cdots<i_r$ be the positions of pivot columns.
            \par Then $\{v_{i_1},v_{i_2},\cdots,v_{i_r}\}$ is a basis for $\operatorname{span}\{v_1,v_2,\cdots,v_n\}$.
            \par Proof see section "Proofs".
        \subsection{Theorem: extend a linearly independent sets to a basis\label{sec:theorem_extension_of_l.i.sets}}
            \hskip \parindent
            \par Let V be a f.d.v.s. over $\mathbb{F}$. Then any l.i. set $\{v_1,v_2,\cdots,v_k\}$of V can be extended to a basis of V.
            \par The proof see section "Proofs".